\chapter{Constraints and Assumptions}
\label{ch:constraints}

\section{Constraints}
\begin{itemize}
\item A smart contract cannot run a machine-learning model. The model therefore runs off-chain, and only its flag and reason reach the contract, through a single authorised oracle address.
\item Development uses a local test chain and test ETH. No real money is involved.
\item Text stored on-chain (flag reasons and appeal notes) is limited to 512 bytes to keep gas costs down.
\item Only open-source tools are used, so the prototype costs nothing to run.
\item The project has one academic year and five team members.
\end{itemize}

\section{Assumptions}
\begin{itemize}
\item The training data, the Kaggle Ethereum Fraud Detection Dataset~\cite{kaggle} (8,731 wallets after removing duplicates), describes wallets and not freelancers. We use the risk of the payout wallet as a proxy until the survey and pilot give us freelancer-history data.
\item The thresholds 0.35 and 0.65 are starting values. They will be tuned on pilot data.
\item Each user has one payout wallet, and payments are made in ETH.
\item The arbiter acts honestly. The contract limits what the arbiter can do (release or refund) but not which of the two they choose.
\item Users have internet access and a modern web browser.
\end{itemize}

\section{Known limitations}
\begin{itemize}
\item The contract has not been audited by an outside party.
\item A single oracle address is a point of trust until several signers are required.
\item Recall of the risk model is 0.61, so some fraudulent wallets will score low. The arbiter and the dispute path are the safeguards.
\end{itemize}
